% Options for packages loaded elsewhere
%DIF LATEXDIFF DIFFERENCE FILE
%DIF DEL ../old/survey.tex   Fri Nov 13 09:45:27 2026
%DIF ADD survey.tex          Fri Nov 13 09:45:27 2026
\PassOptionsToPackage{unicode}{hyperref}
\PassOptionsToPackage{hyphens}{url}
\documentclass[
  12pt,
  a4paper,
]{article}
\usepackage{xcolor}
\usepackage[margin=1in]{geometry}
\usepackage{amsmath,amssymb}
\usepackage[numbers,sort&compress]{natbib}
\bibliographystyle{IEEEtran}
\author{}
\date{}
%DIF PREAMBLE EXTENSION ADDED BY LATEXDIFF
%DIF UNDERLINE PREAMBLE %DIF PREAMBLE
\RequirePackage[normalem]{ulem} %DIF PREAMBLE
\RequirePackage{color}\definecolor{RED}{rgb}{1,0,0}\definecolor{BLUE}{rgb}{0,0,1} %DIF PREAMBLE
\providecommand{\DIFadd}[1]{{\protect\color{blue}\uwave{#1}}} %DIF PREAMBLE
\providecommand{\DIFdel}[1]{{\protect\color{red}\sout{#1}}} %DIF PREAMBLE
%DIF SAFE PREAMBLE %DIF PREAMBLE
\providecommand{\DIFaddbegin}{} %DIF PREAMBLE
\providecommand{\DIFaddend}{} %DIF PREAMBLE
\providecommand{\DIFdelbegin}{} %DIF PREAMBLE
\providecommand{\DIFdelend}{} %DIF PREAMBLE
\providecommand{\DIFmodbegin}{} %DIF PREAMBLE
\providecommand{\DIFmodend}{} %DIF PREAMBLE
%DIF FLOATSAFE PREAMBLE %DIF PREAMBLE
\providecommand{\DIFaddFL}[1]{\DIFadd{#1}} %DIF PREAMBLE
\providecommand{\DIFdelFL}[1]{\DIFdel{#1}} %DIF PREAMBLE
\providecommand{\DIFaddbeginFL}{} %DIF PREAMBLE
\providecommand{\DIFaddendFL}{} %DIF PREAMBLE
\providecommand{\DIFdelbeginFL}{} %DIF PREAMBLE
\providecommand{\DIFdelendFL}{} %DIF PREAMBLE
%DIF AMSMATHULEM PREAMBLE %DIF PREAMBLE
\makeatletter %DIF PREAMBLE
\let\sout@orig\sout %DIF PREAMBLE
\renewcommand{\sout}[1]{\ifmmode\text{\sout@orig{\ensuremath{#1}}}\else\sout@orig{#1}\fi} %DIF PREAMBLE
\makeatother %DIF PREAMBLE
%DIF COLORLISTINGS PREAMBLE %DIF PREAMBLE
\RequirePackage{listings} %DIF PREAMBLE
\RequirePackage{color} %DIF PREAMBLE
\lstdefinelanguage{DIFcode}{ %DIF PREAMBLE
%DIF DIFCODE_UNDERLINE %DIF PREAMBLE
  moredelim=[il][\color{red}\sout]{\%DIF\ <\ }, %DIF PREAMBLE
  moredelim=[il][\color{blue}\uwave]{\%DIF\ >\ } %DIF PREAMBLE
} %DIF PREAMBLE
\lstdefinestyle{DIFverbatimstyle}{ %DIF PREAMBLE
	language=DIFcode, %DIF PREAMBLE
	basicstyle=\ttfamily, %DIF PREAMBLE
	columns=fullflexible, %DIF PREAMBLE
	keepspaces=true %DIF PREAMBLE
} %DIF PREAMBLE
\lstnewenvironment{DIFverbatim}[1][]{\lstset{style=DIFverbatimstyle}}{} %DIF PREAMBLE
\lstnewenvironment{DIFverbatim*}[1][]{\lstset{style=DIFverbatimstyle,showspaces=true}}{} %DIF PREAMBLE
\lstset{extendedchars=\true,inputencoding=utf8}

%DIF END PREAMBLE EXTENSION ADDED BY LATEXDIFF

\begin{document}

\section{Digital Twins in Operation: a Short Survey of the Sample
Corpus}\label{digital-twins-in-operation-a-short-survey-of-the-sample-corpus}

\subsection{1. Introduction}\label{introduction}

The phrase \emph{digital twin} is applied to systems \DIFdelbegin \DIFdel{with very different
levels of }\DIFdelend \DIFaddbegin \DIFadd{of widely differing
}\DIFaddend ambition, and the corpus surveyed here begins by fixing that vocabulary.
A digital model exchanges data with its physical counterpart only
through manual steps \DIFaddbegin \DIFadd{\mbox{%DIFAUXCMD
\citep{sample_ml_anomaly_2023}}\hskip0pt%DIFAUXCMD
}\DIFaddend ; a digital shadow
adds an automatic flow in one direction; a digital twin closes the loop,
with data flowing automatically both ways
\citep{sample_dt_overview_2024}. The distinction is one of data flow
rather than model sophistication, and holding it strictly keeps
expectations honest: in one synthetic survey of forty deployments,
twenty-eight systems described as twins were, by this definition,
shadows \citep{sample_dt_overview_2024}.

This survey groups the sample corpus into three themes -- keeping the
twin true to its asset, using its state stream, and integrating it with
everything else -- and closes with the gaps the corpus leaves open.

\DIFaddbegin \subsection{\DIFadd{4. Integrating with everything
else}}\label{integrating-with-everything-else}

\DIFadd{Twin pilots rarely die of modelling problems; they die of integration.
Interoperability requires agreement at three layers -- transport, syntax
and semantics -- and only the first two can be bought as mature
standards \mbox{%DIFAUXCMD
\citep{sample_std_interop_2021}}\hskip0pt%DIFAUXCMD
. Versioning discipline follows
the same logic; the corpus recommends append-only models, where a field
is deprecated and replaced but never repurposed
\mbox{%DIFAUXCMD
\citep{sample_std_interop_2021}}\hskip0pt%DIFAUXCMD
.
}

\DIFaddend \subsection{3. Using the state stream}\label{using-the-state-stream}

Once a twin exists, its state stream is a naturally clean substrate for
analytics. For anomaly detection, a layered configuration -- a rolling
z-score everywhere, with a learned detector added only where the
false-positive cost justifies its upkeep -- matched an autoencoder's
precision on point faults while remaining immune to the baseline
corruption that made the autoencoder miss the slowest drift entirely
\citep{sample_ml_anomaly_2023}. The cost asymmetry is stark: constant
memory per signal for the baseline against a weekly GPU hour and a human
retraining decision for the learned model
\citep{sample_ml_anomaly_2023}.

\DIFdelbegin \subsection{\DIFdel{4. Integrating with everything
else}}%DIFAUXCMD
\addtocounter{subsection}{-1}%DIFAUXCMD
%DIFDELCMD < \label{integrating-with-everything-else}
%DIFDELCMD < 

%DIFDELCMD < %%%
\DIFdel{Twin pilots rarely die of modelling problems; they die of integration.
Interoperability requires agreement at three layers -- transport, syntax
and semantics -- and only the first two can be bought as mature
standards \mbox{%DIFAUXCMD
\citep{sample_std_interop_2021}}\hskip0pt%DIFAUXCMD
. The practical instrument for
the third is a shared asset-model registry with named owners: sites that
maintained one integrated a new asset in a median of six days, against
seven weeks where semantics were negotiated per project (citation
withdrawn). Versioning discipline follows the same logic; the corpus
recommends append-only models, where a field is deprecated and replaced
but never repurposed \mbox{%DIFAUXCMD
\citep{sample_std_interop_2021}}\hskip0pt%DIFAUXCMD
.
}%DIFDELCMD < 

%DIFDELCMD < %%%
\DIFdelend \subsection{5a. A worked estimate}\label{a.-a-worked-estimate}

The twin's state estimate is corrected by the measurement in proportion
to the gain:

\[
\hat{x}_k = \hat{x}_{k-1} + K_k\,(z_k - \hat{x}_{k-1}), \qquad K_k = \DIFdelbegin \DIFdel{\frac{P_k}{P_k + R}
}\DIFdelend \DIFaddbegin \DIFadd{\frac{P_k}{P_k + R_k}
}\DIFaddend \]


\begin{figure}

{\centering\input{figures/estimator-loop.tex}\par}

\caption{The estimator loop that keeps a twin's state
honest.}\label{fig:estimator-loop}

\end{figure}


\bibliography{/tmp/tmpk2v9x1qe/references.bib}

\end{document}